\section{Related work}\label{sec:related}

\paragraph{LLM agents with validators} Agent frameworks interleave reasoning with actions and use environment
feedback to revise proposals \citep{yao2023react,shinn2023reflexion}. In process systems, LLM agents have been used
for supervisory fault recovery, where proposals are checked against plant knowledge and digital-twin rollouts before
execution \citep{car_todo}. These designs treat validation as mandatory; we ask when it can be skipped and quantify
the safety and compute consequences. \todo{add 2--3 references on LLMs for process control and fault-tolerant control.}

\paragraph{Selective prediction and risk control} Selective classification lets a model abstain when its risk is
high, trading coverage for accuracy \citep{geifman2017selective}. Distribution-free methods calibrate such decisions
to guarantee a risk level with high probability \citep{bates2021rcps,angelopoulos2021ltt}. Our accept rule follows
this tradition: thresholds are chosen on development data so that an upper confidence bound on the failure rate of
accepted proposals stays below a tolerance, and the guarantee is checked once on an independent certification set
with an exact binomial bound \citep{clopper1934}. Unlike standard selective prediction, the screen has three outcomes
(accept, reject, validate) and the cost of a reject is a regeneration by the agent, which we measure in closed loop.

\paragraph{Uncertainty and internal states of LLMs} LLMs' token probabilities carry information about answer
correctness \citep{kadavath2022know}, and linear probes on hidden states detect falsehoods and latent knowledge
\citep{azaria2023internal,burns2023ccs}. Semantic entropy detects confabulations from sampled answers
\citep{farquhar2024semantic}, and probes can approximate it from a single forward pass \citep{kossen2024sep}. These
works ask whether an answer is likely wrong; we ask how many validator calls such signals let an agent skip at a
certified risk, compare them with strong observable baselines, and test them in closed loop.

\paragraph{Attention as evidence of grounding} Whether attention weights explain model behaviour is debated
\citep{jain2019attention,wiegreffe2019attention}; we use attention only as a predictive signal. Lookback Lens
detects contextual hallucinations from the ratio of attention on the context to attention on generated tokens
\citep{chuang2024lookback}. Our region-based grounding is closer to a targeted version of this idea: it measures the
share of attention, while writing each decision-relevant token, that falls on the specific prompt region that decides
that token (an adjacency line in path planning; the relevant plant measurements and knowledge-graph relations in the
CSTR). Our results delimit when such a signal is useful: when failures are failures to retrieve the relevant
information, not failures of reasoning over it.
